%!TeX root=../Memory.tex

\chapter{Introducció}
\label{cap:introduccio}

Aquest capítol presenta el projecte \textbf{Demochain}, un sistema de votació
descentralitzada construït sobre la xarxa \emph{blockchain} d'Algorand i
ancorat periòdicament a Ethereum. L'objectiu d'aquest capítol és situar el
projecte en el seu doble context --acadèmic i tècnic--, exposar la problemàtica
que pretén abordar, repassar els antecedents més rellevants en l'àmbit de la
votació electrònica verificable, justificar la motivació de l'equip i fixar
els objectius i l'abast del que s'entrega. La discussió detallada de cada
aspecte arquitectònic i d'enginyeria es deixa per als capítols posteriors.

%-----------------------------------------------------------------------
\section{Presentació de l'equip}
\label{sec:equip}

El projecte ha estat desenvolupat per un equip de quatre estudiants del
Grau d'Enginyeria Informàtica de la \ac{UIB}. Cap membre de l'equip tenia
experiència prèvia en desenvolupament de \emph{blockchain} ni en
criptografia aplicada en començar el projecte; tots compartim formació en
Python i, en menor grau, en \emph{frontend} amb tecnologies basades en
JavaScript. Aquest punt de partida ha condicionat fortament tant la
planificació --amb una fase inicial de recerca explícita-- com l'abast
final del producte entregat.

\begin{description}
    \item[\textbf{Jordi Vanyó Mestre}] -- Responsable del client web i de la
    integració amb Algorand. Lidera el desenvolupament del \emph{frontend}
    en React i TypeScript, l'arquitectura del repositori i la interacció
    directa amb el contracte d'Algorand (codificació i descodificació
    ARC-4, gestió de \emph{boxes} i signatura de transaccions des del
    navegador).

    \item[\textbf{Dylan Luigi Canning Garcia}] -- Responsable del sistema
    d'ancoratge i del contracte notari a Ethereum. Encarregat del servei en
    Python que llegeix l'estat de la cadena Algorand, en calcula un
    \emph{hash} determinista i el publica al contracte
    \texttt{NotaryContract} desplegat a la \emph{testnet} Sepolia.

    \item[\textbf{Antonio Contestí Coll}] -- Responsable de la coordinació
    del projecte i de la documentació. Encarregat de la gestió del
    \emph{backlog} a GitHub Projects, del seguiment dels \emph{sprints},
    de la coherència de la documentació tècnica i de l'enllaç entre els
    diferents paquets de treball.

    \item[\textbf{Marc Link Cladera}] -- Responsable del \emph{smart
    contract} d'Algorand i de l'arquitectura general. Lidera el disseny
    del contracte \texttt{Demochain} en \emph{algopy}, els diagrames del
    model C4 \cite{brown2018c4} i les decisions arquitectòniques
    recollides als \acp{ADR}.
\end{description}

La distribució de responsabilitats no és estricta. Totes les decisions
arquitectòniques rellevants s'han pres en reunions setmanals conjuntes, i
tot el codi ha entrat a la branca principal a través de \emph{pull
requests} amb revisió creuada (vegeu secció~\ref{sec:versionat} del
capítol~\ref{cap:planificacio}).

%-----------------------------------------------------------------------
\section{Marc acadèmic}
\label{sec:marc-academic}

El projecte s'emmarca en l'assignatura \textbf{21782 -- Laboratori de
Projectes de Software} (Grup 1, Mallorca) del Grau d'Enginyeria
Informàtica de la \acl{UIB} (\ac{UIB}), impartit a l'Escola Politècnica
Superior. Aquesta assignatura té com a propòsit que l'estudiantat
afronti un projecte de programari de complexitat realista aplicant de
manera integrada el conjunt de pràctiques d'enginyeria del programari
adquirides al llarg del grau: anàlisi de requisits, disseny
arquitectònic, planificació iterativa, control de versions, treball en
equip i documentació tècnica. La naturalesa transversal del producte
escollit --un sistema de votació descentralitzat amb dos contractes
intel·ligents, servei d'ancoratge i client web-- ha permès aplicar i
integrar coneixements de diverses matèries: estructures de dades i
algorítmica, sistemes distribuïts, seguretat informàtica, programació
web i bases de dades.

L'execució del projecte s'ha desenvolupat durant el curs acadèmic
\textbf{2025-26}, amb la fase inicial de recerca cap a febrer de 2026 i
l'entrega final el 25 de maig de 2026. Tot el codi font, els
\acp{ADR}, els diagrames C4 i la traça completa del \emph{backlog} estan
disponibles al repositori públic \cite{demochainrepo} sota llicència
oberta, en consonància amb l'esperit acadèmic del projecte.

%-----------------------------------------------------------------------
\section{Context del problema}
\label{sec:context-problema}

La transformació digital de processos col·lectius --consultes,
primàries, eleccions internes d'entitats, referèndums-- ha posat de
manifest una tensió de fons: els sistemes de votació electrònica
desplegats fins ara tendeixen a substituir un punt únic de confiança
(el responsable del recompte físic) per un altre (el responsable del
sistema informàtic). Un servidor central que rep, emmagatzema i
recompta vots és, per definició, una entitat en qui s'han de dipositar
fiançament tant la integritat del recompte com la privacitat del vot
individual. La promesa de la tecnologia \emph{blockchain} aplicada a
la votació és precisament l'eliminació d'aquest punt únic de
confiança.

Els antecedents internacionals d'aquesta problemàtica --casos
desplegats a escala nacional com els d'Estònia, Suïssa o els Estats
Units--, els sistemes acadèmics que defineixen el referent
criptogràfic actual (Helios, Belenios, ElectionGuard) i els projectes
contemporanis sobre \emph{blockchain} (DAVINCI Protocol) es revisen en
detall al capítol~\ref{cap:estat-art} (\emph{Estat de l'art}). En
aquesta introducció ens limitem a constatar que cap d'aquests sistemes
no s'ajusta plenament al cas d'ús d'una organització de mida moderada
--una universitat, una entitat civica o un col·lectiu professional--
que vulgui un sistema obert, descentralitzat, de baix cost operatiu i
amb verificabilitat pública. El nostre projecte explora precisament
aquesta franja.

%-----------------------------------------------------------------------
\section{Motivació}
\label{sec:motivacio}

La motivació del projecte s'articula en tres plànols complementaris:

\subsection{Motivació tècnica}

L'ecosistema \emph{blockchain} ha viscut una maduració significativa
durant els darrers anys. Plataformes com Algorand \cite{chen2017algorand}
ofereixen ara propietats que fan tècnicament viable un sistema de
votació seriós: finalitat de bloc en pocs segons, comissions
menyspreables (de l'ordre de la mil·lèsima de cèntim per transacció), suport
nadiu per a \emph{smart contracts} expressius mitjançant
\emph{algopy} \cite{algorandfoundation2024puya} i operacions
criptogràfiques avançades a nivell de màquina virtual (incloent-hi
opcodes BN254 per a la verificació de proves zk-SNARK). El conjunt
d'aquestes propietats no estava disponible fa cinc anys.

Paral·lelament, Ethereum continua sent la xarxa pública més consolidada
i auditada \cite{buterin2014ethereum}. Tot i el seu cost elevat per a
operacions intensives, la seva fiabilitat com a registre immutable de
fets és superior a la de qualsevol alternativa. Això la fa idònia com a
\emph{capa de notarització} per a fets ocorreguts en altres xarxes.

La combinació d'aquestes dues plataformes --una xarxa permissionada
ràpida i barata per a la lògica electoral, una xarxa pública robusta
per al registre immutable del resultat-- representa una direcció
arquitectònica que mereix exploració empírica.

\subsection{Motivació acadèmica}

El projecte ofereix un escenari pràcticament ideal per posar a prova
de manera coherent un conjunt ampli de pràctiques d'enginyeria del
programari:

\begin{itemize}
    \item Gestió iterativa del \emph{backlog} amb GitHub Projects,
    \emph{sprints} \cite{schwaber2020scrum} i criteris d'acceptació
    en format \emph{Donat/Quan/Aleshores}.
    \item Documentació arquitectònica mitjançant el model C4
    \cite{brown2018c4} i \acp{ADR} per a totes les decisions
    significatives.
    \item Control de versions amb \emph{feature branches}, \emph{pull
    requests} amb revisió creuada obligatòria i \emph{Conventional
    Commits} \cite{conventionalcommits}.
    \item Integració contínua amb GitHub Actions, \emph{git hooks}
    locals i comprovacions automàtiques de format i convenció.
\end{itemize}

La complexitat tècnica del producte (multilingüe quant a
\emph{stack}: Python, TypeScript, Solidity; multinivell: \emph{frontend},
\emph{smart contracts} en dues cadenes, servei d'ancoratge) força que
aquestes pràctiques no siguin un afegit decoratiu sinó una necessitat
operativa.

\subsection{Motivació personal}

A nivell personal, el projecte combina dos ingredients que el fan
atractiu per a un equip d'estudiants en fase final de grau. D'una
banda, l'aprenentatge significatiu d'una tecnologia emergent
--\emph{blockchain}, \emph{smart contracts}-- sobre la qual cap dels
membres tenia coneixement previ. De l'altra, l'oportunitat de treballar
en un problema d'impacte social potencial: la votació és un mecanisme
fonamental de qualsevol col·lectiu organitzat, i contribuir a millorar-ne
la transparència té un valor que va més enllà del purament tècnic.

%-----------------------------------------------------------------------
\section{Objectius del projecte}
\label{sec:objectiu}

L'objectiu general del projecte és \textbf{dissenyar, implementar i
documentar un sistema de votació descentralitzat de propòsit general},
basat en una xarxa permissionada d'Algorand i ancorat a Ethereum, que
serveixi com a prova de concepte tècnicament sòlida sobre la qual es
puguin construir extensions futures --notablement la capa de privacitat
criptogràfica.

Concretament, el sistema entregat ha de satisfer els objectius
específics següents:

\begin{enumerate}
    \item \textbf{Gestió d'organitzacions i cens.} Qualsevol usuari ha de
    poder crear una organització amb nom únic, esdevenir-ne administrador
    automàticament i gestionar un cens d'adreces autoritzades. La càrrega
    massiva d'adreces s'ha de poder fer mitjançant un fitxer \ac{CSV} de
    manera transparent per a l'administrador.

    \item \textbf{Propostes amb fase d'aprovació.} Un membre del cens
    d'una organització ha de poder crear propostes amb un conjunt
    d'opcions i un rang temporal de votació. Abans de la data d'inici,
    la comunitat decideix per quòrum de dos terços si la proposta
    arribarà a celebrar-se com a elecció.

    \item \textbf{Votació preferencial.} Quan una proposta es converteix
    en elecció, els membres del cens han de poder emetre un
    \emph{rànquing} complet de les opcions disponibles. El recompte
    final s'ha de calcular pel mètode de Schulze \cite{schulze2011newmethod},
    que garanteix propietats formals desitjables com la monotonicitat,
    la independència de clons i la consistència amb Condorcet.

    \item \textbf{Verificació pública del vot.} Qualsevol persona,
    sigui o no membre del cens, ha de poder consultar el contingut del
    vot emès per una adreça arbitrària directament des de la cadena
    d'Algorand.

    \item \textbf{Ancoratge multi-node a Ethereum.} En finalitzar una
    elecció, diversos nodes han de llegir independentment l'estat de la
    cadena Algorand, calcular el mateix \emph{hash} determinista i
    enviar-lo al contracte notari d'Ethereum. Quan un nombre suficient
    de nodes coincideix (consens K-de-N), el resultat queda ancorat de
    manera immutable a Sepolia.

    \item \textbf{Reproductibilitat del desplegament.} Tot el sistema
    --\emph{LocalNet} d'Algorand, contracte, \emph{indexer},
    \emph{client} web i serveis auxiliars-- s'ha de poder aixecar
    localment amb una sola comanda mitjançant Docker Compose.

    \item \textbf{Pràctiques d'enginyeria del programari aplicades amb
    rigor.} \emph{Backlog} estructurat en èpiques, històries d'usuari i
    tasques amb criteris d'acceptació explícits; control de versions
    amb \emph{feature branches} i \ac{PR} obligatòria; \emph{Conventional
    Commits}; \acp{ADR} documentats; arquitectura descrita amb el model
    C4.
\end{enumerate}

%-----------------------------------------------------------------------
\section{Contribucions del projecte}
\label{sec:contribucions}

En relació amb els antecedents discutits al capítol~\ref{cap:estat-art},
el projecte \emph{Demochain} aporta una combinació de decisions de
disseny que el diferencien dels sistemes existents:

\begin{itemize}
    \item \textbf{Algorand com a capa primària.} La gran majoria de
    sistemes basats en \emph{blockchain} es construeixen sobre Ethereum
    o variants d'aquesta. La tria d'Algorand com a cadena d'execució del
    contracte de votació és poc habitual i prové d'una anàlisi
    documentada (\ac{ADR} 001) que pondera el cost per transacció, la
    finalitat ràpida del bloc i la disponibilitat d'\emph{smart contracts}
    expressius programables en Python.

    \item \textbf{Ancoratge K-de-N a Ethereum com a notari independent.}
    En lloc de confiar la integritat del resultat exclusivament a la
    cadena d'execució, el projecte hi afegeix una capa explícita de
    notarització: un conjunt de nodes independents calculen el
    \emph{hash} del resultat de manera determinista i el publiquen
    conjuntament al contracte \texttt{NotaryContract} d'Ethereum. La
    propietat aconseguida és que cap entitat individual --ni Algorand,
    ni cap node concret-- pot alterar el resultat unilateralment sense
    que això sigui públicament detectable.

    \item \textbf{Recompte amb el mètode de Schulze.} Els sistemes de
    votació basats en \emph{blockchain} acostumen a oferir recompte per
    pluralitat o, en el millor cas, vot en doble volta. \emph{Demochain}
    implementa el mètode de Schulze \cite{schulze2011newmethod}, que
    permet capturar matisos de preferència que un sistema de
    pluralitat descarta i ofereix propietats formals demostrables.

    \item \textbf{Codi obert pensat per a desplegament real.} El
    repositori inclou el \emph{stack} complet (contracte, client,
    ancoratge, infraestructura amb Docker Compose) llest per ser
    desplegat per una entitat universitària o associativa amb recursos
    modests. No és un experiment de laboratori sinó una base
    productiva.
\end{itemize}

És útil situar el projecte respecte a DAVINCI \cite{vocdoni2026davinci},
que és el sistema contemporani més proper. DAVINCI opta per construir
una capa pròpia de tipus \emph{Layer 2} sobre Ethereum amb una xarxa
de \emph{sequencers} dedicats i un \emph{zkVM} propi. Aquest enfocament
és més ambiciós i tècnicament més pesat, però també exigeix una
infraestructura especialitzada per ser desplegat. \emph{Demochain}
adopta una estratègia complementària: emprar una cadena pública de
propòsit general (Algorand) com a capa de votació, recolzar-se en una
altra cadena pública (Ethereum) com a capa de notarització, i mantenir
el component d'ancoratge com un servei petit i auditables en Python
que qualsevol universitat o entitat pot operar amb recursos modestos.
Es tracta, doncs, de dues respostes diferents a una mateixa pregunta
de fons: com aconseguir verificabilitat pública i resistència a
manipulació sense centralitzar la confiança.

Cal subratllar, finalment, que aquestes contribucions són de disseny i
d'integració. El projecte no aporta cap primitiva criptogràfica
original ni cap protocol nou: tots els components criptogràfics emprats
estan ben establerts a la literatura.

%-----------------------------------------------------------------------
\section{Abast i no objectius}
\label{sec:abast}

És important deixar constància, ja a la introducció, dels límits del
que s'entrega. El \ac{MVP} actual \textbf{no inclou} privacitat
criptogràfica del vot. Concretament, queden fora de l'abast d'aquesta
versió:

\begin{itemize}
    \item Xifratge homomòrfic dels vots (esquema ElGamal exponencial), tal
    com el descriu Adida \cite{adida2008helios}.
    \item Proves de coneixement zero (\acsp{zkSNARK} \emph{Groth16}) per
    acreditar l'elegibilitat del votant sense revelar-ne la identitat.
    \item Generació distribuïda de claus i cerimònia de desxifrat amb
    \emph{trustees} i proves Chaum-Pedersen, segons l'esquema de
    Belenios \cite{cortier2019belenios}.
    \item Integració amb un proveïdor d'identitat (\ac{IdP}) extern per a
    l'autenticació del votant.
\end{itemize}

Aquests components estaven contemplats al \emph{backlog} inicial
(èpiques E3, E4, E5 i E6 al repositori \cite{demochainrepo}) i es van
descartar de l'abast del \ac{MVP} al final del primer mes de
desenvolupament, en constatar que la seva complexitat depassava
l'horitzó temporal del projecte. La justificació detallada d'aquesta
decisió i la seva planificació com a treball futur es recullen al
capítol~\ref{cap:planificacio} (secció~\ref{sec:reduccio-abast}) i al
capítol de conclusions.

Cal subratllar que el sistema entregat \textbf{sí ofereix
verificabilitat pública} del resultat: qualsevol persona pot consultar
els vots individuals a la cadena d'Algorand i comprovar que el
\emph{hash} ancorat a Ethereum coincideix amb l'estat real de
l'elecció. Aquesta propietat ja és, per ella mateixa, una millora
substancial respecte als sistemes de votació centralitzats
convencionals; el que queda fora és la capa de \textbf{privacitat}
criptogràfica, que protegiria el vot individual de la inspecció
pública.

%-----------------------------------------------------------------------
\section{Enfocament metodològic}
\label{sec:metodologia}

L'execució del projecte ha adoptat un enfocament \textbf{iteratiu i
incremental} inspirat en Scrum \cite{schwaber2020scrum}, amb les
adaptacions necessàries per a un equip acadèmic petit. La metodologia
es resumeix en quatre pràctiques principals que es desenvolupen amb
detall al capítol~\ref{cap:planificacio}:

\begin{enumerate}
    \item \textbf{Recerca prèvia documentada.} Abans de prendre cap
    decisió arquitectònica significativa, s'ha investigat l'espai de
    solucions i s'ha documentat la decisió en un \acl{ADR} (\ac{ADR}).
    Els dos \acp{ADR} actuals (\acl{ADR} 000 sobre llenguatge de
    \emph{smart contracts} i \acl{ADR} 001 sobre plataforma
    \emph{blockchain}) consten al repositori.

    \item \textbf{Planificació amb \emph{backlog} jerarquitzat.}
    Èpiques, històries d'usuari, \emph{spikes} i tasques estructurats
    a GitHub Projects amb identificadors traçables (E1, E2-US1,
    E2-US1-T3, etc.) i criteris d'acceptació formals.

    \item \textbf{Arquitectura documentada amb C4.} Diagrames de
    context, contenidors i components mantinguts al directori
    \texttt{docs/diagrams/} del repositori, ampliats a mesura que el
    sistema creix \cite{brown2018c4}.

    \item \textbf{Control de versions amb revisió creuada
    obligatòria.} Tot el codi entra a la branca principal via
    \emph{pull request}, amb \emph{check} automàtic de format i
    convenció i revisió manual d'almenys un altre membre de l'equip.
\end{enumerate}

%-----------------------------------------------------------------------
\section{Estructura del document}
\label{sec:estructura}

La resta del document s'organitza en sis capítols més, amb una
estructura que combina les fases conceptuals d'un projecte d'enginyeria
del programari amb una secció final de tancament.

\begin{itemize}
    \item El \textbf{capítol~\ref{cap:estat-art} (Estat de l'art)}
    revisa les experiències nacionals de votació electrònica
    desplegades a Estònia, Suïssa i els Estats Units, els sistemes
    acadèmics consolidats (Helios, Belenios, ElectionGuard) i els
    projectes contemporanis sobre \emph{blockchain} (DAVINCI
    Protocol), i situa el nostre projecte en aquest panorama.

    \item El \textbf{capítol~3 (Anàlisi)} identifica les parts
    interessades del sistema, exposa l'especificació funcional --casos
    d'ús, històries d'usuari i requisits de producte-- i
    l'especificació no funcional --requisits de qualitat i restriccions
    tecnològiques, legals i de calendari. Inclou també el llenguatge
    ubic del domini, el model de dades de les estructures emmagatzemades
    al contracte i una síntesi de la recerca prèvia sobre
    \emph{blockchain} i criptografia.

    \item El \textbf{capítol~4 (Disseny)} presenta l'arquitectura del
    sistema mitjançant el model C4 (diagrames de context, contenidors
    i components per a cada peça del sistema), descriu el patró de
    desplegament basat en Docker Compose i exposa el disseny de la
    interfície d'usuari amb les pantalles del client i els diagrames
    de seqüència de les interaccions principals.

    \item El \textbf{capítol~\ref{cap:planificacio} (Planificació)}
    detalla la gestió del cicle de vida del projecte: l'estructura del
    \emph{backlog}, les \aclu{DoR} (\ac{DoR}) i \aclu{DoD} (\ac{DoD}),
    els \emph{sprints} executats, els camins crítics identificats i la
    política de versionat del codi.

    \item El \textbf{capítol~6 (Desenvolupament)} descriu l'execució
    real del projecte: l'organització de l'equip, la traça de les
    iteracions, les versions generades, els problemes trobats i les
    seves resolucions, i les demostracions funcionals del producte
    final.

    \item El \textbf{capítol~7 (Conclusions)} valora el producte
    entregat respecte als objectius inicials, exposa el que s'ha fet,
    el que no s'ha fet i el que queda per fer, reflexiona sobre el
    funcionament de l'equip i les lliçons apreses, i identifica les
    línies de treball futures més rellevants.
\end{itemize}

El codi font del projecte, els \acp{ADR}, els diagrames C4 i la traça
completa del \emph{backlog} estan disponibles al repositori públic de
GitHub \cite{demochainrepo}.
